% Source: https://artofproblemsolving.com/community/c3h1543718_probability_problem
% Stable identity: aops-topic-1543718-q01
PHYSICS uses P, H, Y, C once each, S twice and I once. The bag has two S tiles and two I tiles. Treat the $12$ tiles as distinguishable.

\textbf{(a)} All $\dbinom{12}{7}=792$ sets of $7$ tiles are equally likely. A favourable set contains P, H, Y, C, both S tiles and one of the two I tiles: $2$ sets. The probability is $\dfrac{2}{792}=\boxed{\dfrac{1}{396}}$.

\textbf{(b)} All $12\cdot11\cdot10\cdot9\cdot8\cdot7\cdot6=3991680$ ordered draws are equally likely. To spell P, H, Y, S, I, C, S in order there are $2$ choices for the first S tile (the second S is then forced) and $2$ choices for the I tile: $4$ favourable draws. The probability is
\[
\dfrac{4}{3991680}=\boxed{\dfrac{1}{997920}} .
\]
(Consistently, (b) equals (a) times the probability $\dfrac{2!}{7!}=\dfrac{1}{2520}$ that $7$ tiles forming the right set come out in a correct order, the two S tiles being interchangeable.)
